%------------------------------------------------------------------------------%
\section{Volume pelo método das cascas cilíndricas}
\label{sec:Volume_cascas_cilindricas}

%------------------------------------------------------------------------------%
\begin{figure}
  \centering
  \includegraphics[width=0.5\linewidth]{volume-cascas-01}
  \caption{Volume por cascas cilíndricas -- 01}
  \label{fig:volume_cascas_01}
\end{figure}
%------------------------------------------------------------------------------%

%------------------------------------------------------------------------------%
\begin{figure}
  \centering
  \includegraphics[width=0.5\linewidth]{volume-cascas-02}
  \caption{Volume por cascas cilíndricas -- 02}
  \label{fig:volume_cascas_02}
\end{figure}
%------------------------------------------------------------------------------%

%------------------------------------------------------------------------------%
\begin{figure}
  \centering
  \includegraphics[width=0.5\linewidth]{volume-cascas-03}
  \caption{Volume por cascas cilíndricas -- 03}
  \label{fig:volume_cascas_03}
\end{figure}
%------------------------------------------------------------------------------%

%------------------------------------------------------------------------------%
\begin{figure}
  \centering
  \includegraphics[width=0.5\linewidth]{volume-cascas-04}
  \caption{Volume por cascas cilíndricas -- 04}
  \label{fig:volume_cascas_04}
\end{figure}
%------------------------------------------------------------------------------%

%------------------------------------------------------------------------------%
\begin{figure}
  \centering
  \includegraphics[width=0.5\linewidth]{volume-cascas-05}
  \caption{Volume por cascas cilíndricas -- 05}
  \label{fig:volume_cascas_05}
\end{figure}
%------------------------------------------------------------------------------%

%------------------------------------------------------------------------------%
\begin{figure}
  \centering
  \includegraphics[width=0.5\linewidth]{volume-cascas-06}
  \caption{Volume por cascas cilíndricas -- 06}
  \label{fig:volume_cascas_06}
\end{figure}
%------------------------------------------------------------------------------%

%------------------------------------------------------------------------------%
\begin{figure}
  \centering
  \includegraphics[width=0.5\linewidth]{volume-cascas-07}
  \caption{Volume por cascas cilíndricas -- 07}
  \label{fig:volume_cascas_07}
\end{figure}
%------------------------------------------------------------------------------%

%------------------------------------------------------------------------------%
\begin{figure}
  \centering
  \includegraphics[width=0.5\linewidth]{volume-cascas-08}
  \caption{Volume por cascas cilíndricas -- 08}
  \label{fig:volume_cascas_08}
\end{figure}
%------------------------------------------------------------------------------%

%------------------------------------------------------------------------------%
\begin{figure}
  \centering
  \includegraphics[width=0.5\linewidth]{volume-cascas-09}
  \caption{Volume por cascas cilíndricas -- 09}
  \label{fig:volume_cascas_09}
\end{figure}
%------------------------------------------------------------------------------%

%------------------------------------------------------------------------------%
\begin{definicao}{Volume por cascas cilíndricas}{def:volume-cascas}
  O volume do sólido $S$ obtido pela rotação em torno do eixo $y$ da região
  \[
    R = \big\{(x, y)\colon 0 \leq y \leq f(x), a \leq x \leq b \big\}
  \]
  é calculado pela integral
  \begin{align*}
    V & = \int_a^b 2\pi\,(\text{raio da casca})\,(\text{altura da casca})\, dx \\
      & = \int_a^b 2\pi x f(x) dx
  \end{align*}
\end{definicao}
%------------------------------------------------------------------------------%

%------------------------------------------------------------------------------%
\begin{example}
  Encontre o volume do sólido obtido pela rotação em torno do eixo
  $y$ da região delimitada por $y = 2x^2 - x
  ^3$ e $y = 0$.
\end{example}
%------------------------------------------------------------------------------%

%------------------------------------------------------------------------------%
\begin{example}
  Encontre o volume do sólido obtido pela rotação em torno do eixo
  $y$ da região entre $y = x$ e $y = x^2$.
\end{example}
%------------------------------------------------------------------------------%

%------------------------------------------------------------------------------%
\begin{example}
  Encontre o volume do sólido obtido pela rotação da região
  delimitada por $y = x - x^2$ e $y = 0$ em torno da reta $x = 2$.
\end{example}
%------------------------------------------------------------------------------%

Ao calcular volumes por cascas cilíndricas fique atento aos pontos:
\begin{enumerate}
  \item Se o sólido é obtido pela rotação de uma região em torno do eixo $x$
    (ou em torno de uma reta paralela ao eixo $x$) o raio da casca e a altura
    são funções de $y$ e o volume é dado por uma integral na variável $y$.

  \item Se o sólido é obtido pela rotação de uma região em torno do eixo $y$
    (ou em torno de uma reta paralela ao eixo $y$) o raio da casca e a altura
    são funções de $x$ e o volume é dado por uma integral na variável $x$.
\end{enumerate}

%------------------------------------------------------------------------------%
